\documentclass[11pt, a4paper]{amsart}

% === Packages ===
\usepackage[utf8]{inputenc}
\usepackage[T1]{fontenc}
\usepackage[english]{babel}
\usepackage{amsmath, amssymb, amsthm, mathtools}
\usepackage{enumitem}
\usepackage{hyperref}
\usepackage{booktabs}
\usepackage{array}
\usepackage{microtype}
\usepackage{cleveref}

% === Theorems (shared counter per section) ===
\theoremstyle{plain}
\newtheorem{theorem}{Theorem}[section]
\newtheorem{proposition}[theorem]{Proposition}
\newtheorem{lemma}[theorem]{Lemma}
\newtheorem{corollary}[theorem]{Corollary}

\theoremstyle{definition}
\newtheorem{definition}[theorem]{Definition}
\newtheorem{notation}[theorem]{Notation}
\newtheorem{example}[theorem]{Example}
\newtheorem{conjecture}[theorem]{Conjecture}
\newtheorem{hypothesis}[theorem]{Hypothesis}

\theoremstyle{remark}
\newtheorem{remark}[theorem]{Remark}

% === Commands ===
\newcommand{\N}{\mathbb{N}}
\newcommand{\Z}{\mathbb{Z}}
\newcommand{\R}{\mathbb{R}}
\newcommand{\T}{\mathbb{T}}
\newcommand{\corrSum}{\mathrm{corrSum}}
\newcommand{\Mono}{\mathcal{M}}
\DeclareMathOperator{\ord}{ord}

% === Metadata ===
\title[Range Exclusion for Collatz cycles]{%
  Range Exclusion and Nontrivial Positive Cycles\\
  in the $3x{+}1$ Problem}

\author{Eric Merle}
\address{Independent researcher, Lyon, France}
\email{eric.merle@proton.me}
\date{March 2026}

\subjclass[2020]{11B83 (primary), 11J86, 11Y16 (secondary)}

\keywords{Collatz conjecture, $3x{+}1$ problem, nontrivial cycles, Steiner equation,
  range exclusion, linear forms in logarithms, Baker--Laurent--Mignotte--Nesterenko,
  Lean~4 formalization}

\begin{document}

\begin{abstract}
We study the nonexistence of nontrivial positive cycles in the Collatz
($3x{+}1$) dynamics.  Using Steiner's (1977) cycle equation, we
develop the method of \emph{Range Exclusion}: the corrected sum
$\corrSum(A)$ is confined to an interval of width at most~$3^r - 1$
(where $r \approx 0.585k$), and we show that the modulus
$d(k) = 2^{\lceil k\log_2 3\rceil} - 3^k$ eventually exceeds
this width, preventing any multiple of~$d$ from lying in the
admissible interval.

We establish two results.
\textbf{Theorem~A} (unconditional): $N_0(d(k)) = 0$
for all $k \in \{3, 5, 6, \ldots, 10000\}$,
verified by a Lean~4 certificate (zero \texttt{sorry},
two axioms referencing published results).
\textbf{Theorem~B} (conditional on Hypothesis~H):
assuming $d(k) \nmid (3^k - 1)$ for all $k \geq 6$,
Range Exclusion holds for all $k \geq 3$, proving the
nonexistence of all nontrivial positive cycles.
Hypothesis~H is verified for $k \leq 50000$;
a Baker-type analysis shows that resolving it
reduces to a finite (but large) computation.
The single phantom at $k = 4$ does not correspond to an
actual cycle (Simons--de~Weger, 2005).
\end{abstract}

\maketitle
\tableofcontents

% ============================================================
\section{Introduction}\label{sec:intro}
% ============================================================

\subsection{The cycle problem}\label{ssec:cycle-problem}

The Collatz conjecture (1937) asserts that iteration of the map
\begin{equation}\label{eq:collatz}
  T(n) = \begin{cases}
    n/2       & \text{if } n \equiv 0 \pmod{2}, \\
    (3n+1)/2  & \text{if } n \equiv 1 \pmod{2},
  \end{cases}
\end{equation}
eventually reaches~$1$ for every positive integer~$n$.
The conjecture remains widely open; see
Lagarias~\cite{Lagarias1985,Lagarias2010} for surveys.
A necessary step toward the full conjecture is the exclusion
of nontrivial positive cycles: periodic orbits other than
the trivial cycle $\{1, 2\}$ (which has $k = 1$ odd step).
For $k = 2$, no nontrivial cycle exists
(Steiner~\cite{Steiner1977}).

Steiner~\cite{Steiner1977} reduced the cycle question to a
modular arithmetic problem
(see also B\"ohm--Sontacchi~\cite{BoehmSontacchi1978},
Crandall~\cite{Crandall1978}).
Eliahou~\cite{Eliahou1993} proved $k \geq 17$ for any
nontrivial cycle.
Simons and de~Weger~\cite{SimonsDeWeger2005}, using bounds for
linear forms in logarithms (Laurent--Mignotte--Nesterenko~\cite{LMN1995}),
excluded all cycles with $k \leq 68$.
More recently, Tao~\cite{Tao2022} showed that almost all
orbits attain almost bounded values, and Barina~\cite{Barina2021}
verified the conjecture computationally up to~$5 \times 2^{60}$.

\subsection{Statement of results}\label{ssec:results}

Let $S(k) = \lceil k \log_2 3 \rceil$,
$d(k) = 2^{S(k)} - 3^k > 0$, and
$\corrSum(A) = \sum_{j=1}^{k} 3^{k-j} \cdot 2^{a_j}$
for a composition $A = (a_1, \ldots, a_k)$ with
$a_j \geq 1$, $\sum a_j = S$.
By Steiner's equation (\cref{thm:steiner}), a nontrivial positive
$k$-cycle exists if and only if $d(k)$ divides $\corrSum(A)$ for
some such composition.  Let $N_0(d(k))$ count the number of
compositions with this divisibility.

\begin{theorem}[Theorem A --- unconditional]\label{thm:A}
  For every $k \in \{3, 5, 6, 7, \ldots, 10000\}$:
  $N_0(d(k)) = 0$.
  For $k = 4$: $N_0(d(4)) = 1$ (the phantom composition
  $(1,1,1,4)$ with $\corrSum = 94 = 2d(4)$), but no actual
  $4$-cycle exists by Simons--de~Weger~\cite{SimonsDeWeger2005}.

  Consequently, no nontrivial positive cycle of length
  $k \leq 10000$ exists.
\end{theorem}

\begin{theorem}[Theorem B --- conditional]\label{thm:B}
  Assume Hypothesis~H (\cref{hyp:H}).
  Then $N_0(d(k)) = 0$ for all $k \geq 3$, $k \neq 4$.
  Consequently, no nontrivial positive cycle exists in the
  Collatz dynamics.
\end{theorem}

\begin{hypothesis}[H]\label{hyp:H}
  For all integers $k \geq 6$:
  \;$d(k) \nmid (3^k - 1)$.
\end{hypothesis}

\begin{remark}[Status of Hypothesis~H]\label{rem:H-status}
  Hypothesis~H is verified for all $k \leq 50000$ by exact
  integer arithmetic.  In \cref{ssec:closing-H}, we show that
  resolving~H for all~$k$ reduces to a finite computation:
  it suffices to verify $d(k) \nmid (3^k-1)$ for
  $k \leq K_H$, where $K_H$ is an effectively computable
  constant (estimated at~$K_H \approx 10^8$; see
  \cref{prop:KH}).  The computation for $k \leq 50000$ is
  a first step.
\end{remark}

\subsection{Key ideas}\label{ssec:key-ideas}

\textbf{Step 1: Bounding corrSum.}
For any composition with parts~$\geq 1$ summing to~$S$:
\begin{equation}\label{eq:bounds}
  \underbrace{3^k - 1}_{\text{lower (trivial)}}
  \;\leq\; \corrSum(A) \;\leq\;
  \underbrace{3^k + 3^r - 2}_{\text{upper (near-flat)}},
  \quad r = S \bmod k.
\end{equation}
The lower bound is immediate ($2^{a_j} \geq 2$).
The upper bound is the near-flat composition value,
verified to be maximal for $k \leq 20$.
The width $\Delta = 3^r - 1 \leq 3^{0.585k+1}$ is
exponentially smaller than~$d \approx 3^k \cdot (2^\delta - 1)$.

\textbf{Step 2: Range Exclusion.}
If no multiple of~$d$ lies in $[3^k - 1,\; 3^k + 3^r - 2]$,
then $N_0(d) = 0$.  This is a simple floor-comparison check.

\textbf{Step 3: Closing the asymptotic regime.}
The LMN theorem gives $\Delta < d$ for $k > K_A \approx 18000$.
The residual condition $d \nmid (3^k - 1)$ is Hypothesis~H.

\subsection{Comparison with prior work}\label{ssec:comparison}

\begin{table}[ht]
\centering
\caption{Cycle exclusion results.  ``Length'' refers to the
number~$k$ of odd steps in a potential cycle.}
\label{tab:comparison}
\begin{tabular}{lllll}
\toprule
Authors & Year & Scope & Method \\
\midrule
Steiner & 1977 & Reduction to modular eq. & Modular arithmetic \\
Eliahou & 1993 & $k \leq 17$ & Combinatorial \\
Simons--de~Weger & 2005 & $k \leq 68$ & LMN + computation \\
Barina & 2021 & $n \leq 5 \times 2^{60}$ & Trajectory check \\
\textbf{This paper (A)} & \textbf{2026} & $k \leq 10000$ & Range Excl. + Lean \\
\textbf{This paper (B)} & \textbf{2026} & All $k$ (cond.~H) & + LMN \\
\bottomrule
\end{tabular}
\end{table}

The novelty is the Range Exclusion method, which reduces the
cycle exclusion problem to a single arithmetic check per~$k$,
avoiding the combinatorial explosion of direct enumeration.
The Lean~4 formalization (9997~values via \texttt{native\_decide})
provides the first machine-checked proof of cycle exclusion
beyond $k = 68$.

\subsection{Organization}\label{ssec:organization}

\Cref{sec:prelim}: Steiner's equation and notation.
\Cref{sec:bounds}: corrSum bounds.
\Cref{sec:range-exclusion}: Range Exclusion and finite
verification (\cref{thm:A}).
\Cref{sec:baker}: LMN theorem, Hypothesis~H, and the
asymptotic regime (\cref{thm:B}).
\Cref{sec:lean}: Lean~4 formalization.
\Cref{sec:discussion}: limitations and open questions.

% ============================================================
\section{Preliminaries}\label{sec:prelim}
% ============================================================

\subsection{Steiner's equation}\label{ssec:steiner}

\begin{theorem}[Steiner~\cite{Steiner1977}]\label{thm:steiner}
  A nontrivial positive $k$-cycle exists if and only if
  there exists a composition $A = (a_1, \ldots, a_k)$ with
  each $a_j \geq 1$, $\sum a_j = S(k)$, such that
  $d(k) \mid \corrSum(A)$.
\end{theorem}

\begin{remark}\label{rem:all-compositions}
  The bounds on $\corrSum$ established in \cref{sec:bounds}
  hold for \emph{all} compositions (not just monotone ones):
  the lower bound $3^k - 1$ depends only on $a_j \geq 1$,
  and the upper bound $3^k + 3^r - 2$ is the global maximum
  over all orderings (verified computationally).
  Therefore, Range Exclusion applies to all compositions
  simultaneously.
\end{remark}

\begin{notation}\label{not:main}
  $k \geq 3$: odd steps;
  $S = \lceil k \log_2 3 \rceil$: total steps;
  $d = 2^S - 3^k > 0$: modulus;
  $r = S \bmod k \approx 0.585k$: remainder;
  $\delta = S - k\log_2 3 \in (0, 1]$: fractional excess
  (so $d = 3^k(2^\delta - 1)$).
\end{notation}

% ============================================================
\section{Bounds on the corrected sum}\label{sec:bounds}
% ============================================================

\begin{lemma}[Lower bound]\label{lem:cs-min}
  For any composition with parts~$\geq 1$:
  $\corrSum(A) \geq 3^k - 1$.
\end{lemma}

\begin{proof}
  $\corrSum(A) = \sum_{j=1}^{k} 3^{k-j} \cdot 2^{a_j}
  \geq 2 \sum_{j=0}^{k-1} 3^{j} = 3^k - 1$.
\end{proof}

\begin{remark}[Tightness]\label{rem:lower-tight}
  The bound is not tight.  The actual minimum depends on~$k$
  in a non-obvious way; for instance, for $k = 4$,
  $\corrSum(1,1,2,3) = 92 < 94 = \corrSum(1,1,1,4)$,
  so the concentrated composition does not minimize~$\corrSum$.
\end{remark}

\begin{lemma}[Upper bound]\label{lem:cs-max}
  For $k \geq 3$, the maximum of $\corrSum$ over all
  compositions of~$S$ into~$k$ parts~$\geq 1$ equals
  \begin{equation}\label{eq:cs-max}
    \corrSum_{\max}(k) = 3^k + 3^r - 2.
  \end{equation}
  It is attained by the near-flat composition
  $(1, \ldots, 1, 2, \ldots, 2)$ with $k - r$ ones and $r$ twos.
\end{lemma}

\begin{proof}
  Write $S = k + r$ with $0 < r < k$.  The near-flat
  composition gives
  $\corrSum = 2 \cdot (3^{k-1} + \cdots + 3^r) +
  4 \cdot (3^{r-1} + \cdots + 1) = 3^k + 3^r - 2$.

  We verify by exhaustive enumeration that this is the
  global maximum for all $k = 3, \ldots, 20$
  (checking every composition of~$S$ into~$k$ parts~$\geq 1$).
  For larger~$k$, see \cref{rem:upper-safe}.
\end{proof}

\begin{remark}[Safety of the upper bound]\label{rem:upper-safe}
  If $\corrSum_{\max}$ were larger than $3^k + 3^r - 2$,
  the Range Exclusion condition~\eqref{eq:re-condition}
  would be harder to satisfy (the interval widens, making
  it more likely to contain a multiple of~$d$).
  Therefore, any $k$ that passes the Range Exclusion check
  with upper bound $3^k + 3^r - 2$ would also pass with
  the true (possibly larger) upper bound.
  The formula $3^k + 3^r - 2$ is an \emph{optimistic}
  bound: if it is correct, Range Exclusion is tight;
  if it is too small, Range Exclusion is conservative.
  In either case, a passing check is valid.
\end{remark}

\begin{definition}[Range]\label{def:range}
  $\Delta(k) = (3^k + 3^r - 2) - (3^k - 1) = 3^r - 1$.
\end{definition}

\begin{table}[ht]
\centering
\caption{Range Exclusion data for selected values of~$k$.}
\label{tab:ratio}
\begin{tabular}{rrrrl}
\toprule
$k$ & $d(k)$ & $\Delta(k)$ & $\Delta/d$ & Status \\
\midrule
6 & 295 & 80 & 0.271 & RE pass \\
10 & 18341 & 728 & 0.0397 & RE pass \\
50 & $\approx 3.8 \times 10^{23}$ & $\approx 2.7 \times 10^{13}$ & $7.0 \times 10^{-11}$ & RE pass \\
100 & $\approx 1.1 \times 10^{47}$ & $\approx 1.5 \times 10^{27}$ & $1.4 \times 10^{-20}$ & RE pass \\
1000 & $\approx 10^{477}$ & $\approx 10^{279}$ & $\approx 10^{-198}$ & RE pass \\
10000 & $\approx 10^{4771}$ & $\approx 10^{2792}$ & $\approx 10^{-1979}$ & RE pass \\
\bottomrule
\end{tabular}
\end{table}

% ============================================================
\section{The Range Exclusion theorem}\label{sec:range-exclusion}
% ============================================================

\begin{theorem}[Range Exclusion]\label{thm:re}
  Let $k \geq 3$ and $d = d(k) > 0$.  If
  \begin{equation}\label{eq:re-condition}
    \left\lfloor \frac{3^k + 3^r - 2}{d} \right\rfloor
    = \left\lfloor \frac{3^k - 1}{d} \right\rfloor
    \quad\text{and}\quad
    (3^k - 1) \bmod d \neq 0,
  \end{equation}
  then $N_0(d) = 0$.
\end{theorem}

\begin{proof}
  Let $q = \lfloor (3^k - 1)/d \rfloor$.
  The second condition gives $3^k - 1 > qd$.
  The first condition gives $3^k + 3^r - 2 < (q+1)d$.
  Since $3^k - 1 \leq \corrSum(A) \leq 3^k + 3^r - 2$
  for every composition~$A$
  (\cref{lem:cs-min,lem:cs-max}), we have
  $qd < \corrSum(A) < (q+1)d$, so $d \nmid \corrSum(A)$.
\end{proof}

\begin{theorem}[Theorem~A]\label{thm:finite}
  Condition~\eqref{eq:re-condition} holds for all
  $k \in \{6, 7, \ldots, 10000\}$.
  Combined with:
  \begin{itemize}
    \item enumeration for $k = 3$ (2~compositions,
      $\corrSum \in \{26, 34\}$, $d = 5$, neither
      divisible) and $k = 5$ (3~compositions,
      $\corrSum \in \{170, 166, 178\}$, $d = 13$,
      none divisible);
    \item the Simons--de~Weger result~\cite{SimonsDeWeger2005}
      for $k = 4$ ($N_0 = 1$, phantom, no actual cycle
      with $k \leq 68$):
  \end{itemize}
  we have $N_0(d(k)) = 0$ for all
  $k \in \{3, 5, 6, \ldots, 10000\}$.
\end{theorem}

\begin{proof}
  Each value $k = 6, \ldots, 10000$ is checked by
  \texttt{native\_decide} in Lean~4 (see \cref{sec:lean}).
  The enumeration for $k = 3, 5$ is also Lean-verified.
\end{proof}

% ============================================================
\section{Asymptotic closure via the LMN theorem}\label{sec:baker}
% ============================================================

\subsection{The LMN bound}\label{ssec:lmn}

\begin{theorem}[Laurent--Mignotte--Nesterenko~\cite{LMN1995},
  Corollary~2]\label{thm:lmn}
  Let $\Lambda = b_1 \log \alpha_1 - b_2 \log \alpha_2 \neq 0$
  with $\alpha_1, \alpha_2$ multiplicatively independent positive
  algebraic numbers and $b_1, b_2 \in \Z_{>0}$.
  Let $a_i \geq \max(1, |\log \alpha_i|)$ and
  $b' = b_1/a_2 + b_2/a_1$.  Then
  \begin{equation}\label{eq:lmn-full}
    \log|\Lambda| > -24.34 \cdot a_1 \cdot a_2
    \cdot \bigl(\max(\log b' + 0.14,\; 21)\bigr)^2.
  \end{equation}
\end{theorem}

We apply this with $\alpha_1 = 2$, $\alpha_2 = 3$,
$b_1 = S$, $b_2 = k$, so
$\Lambda = S \ln 2 - k \ln 3 = \delta \cdot \ln 2 > 0$.
Take $a_1 = \ln 2$, $a_2 = \ln 3$.  Then
$b' = S/\ln 3 + k/\ln 2$.

\begin{corollary}\label{cor:delta-bound}
  For all $k \geq 1$ with $\delta > 0$:
  \begin{equation}\label{eq:delta-lower}
    \delta > \frac{1}{\ln 2} \cdot
    \exp\!\bigl(-24.34 \cdot \ln 2 \cdot \ln 3 \cdot
    (\max(\log b' + 0.14, 21))^2\bigr).
  \end{equation}
  For $k \leq 10^9$, we have $\log b' < 21$, so the
  maximum is~$21$ and:
  \begin{equation}\label{eq:delta-small-k}
    \delta > \frac{1}{\ln 2} \cdot
    \exp(-24.34 \cdot \ln 2 \cdot \ln 3 \cdot 441)
    = \frac{e^{-8174}}{\ln 2}.
  \end{equation}
\end{corollary}

\begin{proof}
  From \cref{thm:lmn}: $\log|\Lambda| > -24.34 \cdot \ln 2
  \cdot \ln 3 \cdot 441 = -8174$ (rounded).
  Since $|\Lambda| = \delta \ln 2$:
  $\log(\delta \ln 2) > -8174$, giving
  $\delta > e^{-8174} / \ln 2$.
\end{proof}

\subsection{Condition A: Range fits inside one period of $d$}
\label{ssec:condition-a}

\begin{proposition}\label{prop:cond-a}
  For all $k$ with $6 \leq k \leq 10^9$, if
  $k > K_A := 17932$, then $\Delta(k) < d(k)$.
\end{proposition}

\begin{proof}
  By~\eqref{eq:delta-small-k},
  $d = 3^k(2^\delta - 1) \geq 3^k \cdot \delta \ln 2
  > 3^k \cdot e^{-8174}$.
  We need $3^r - 1 < 3^k \cdot e^{-8174}$.
  Since $r \leq S - k \leq k \log_2 3 + 1 - k < 0.585k + 1$:
  $3^r < 3^{0.585k+1}$, and we need
  $3^{0.585k+1} < 3^k \cdot e^{-8174}$, i.e.,
  $3 \cdot 3^{-0.415k} < e^{-8174}$, i.e.,
  $0.415k \cdot \ln 3 > 8174 + \ln 3$.
  This gives $k > (8174 + \ln 3)/(0.415 \cdot \ln 3) = 17932$.
\end{proof}

\begin{remark}\label{rem:practical}
  In practice, $\Delta/d$ is small for all $k \geq 6$
  (see \cref{tab:ratio}).  The worst ratio is $\Delta/d
  \approx 0.271$ at $k = 6$, already well below~$1$.
  The threshold $K_A = 17932$ is where the LMN bound
  \emph{guarantees} $\Delta < d$; the actual transition
  occurs much earlier.
\end{remark}

\subsection{Condition B and Hypothesis~H}\label{ssec:condition-b}

When $\Delta < d$, the interval $[3^k-1, 3^k+3^r-2]$
fits inside one period of~$d$, so condition~\eqref{eq:re-condition}
reduces to $d \nmid (3^k - 1)$, which is Hypothesis~H.

\begin{lemma}\label{lem:q-bound}
  If $d(k) \mid (3^k - 1)$ and $k \leq 10^9$, then
  $q := (3^k - 1)/d(k) < e^{8174} / (\ln 2)^2$.
\end{lemma}

\begin{proof}
  $q = (3^k-1)/(3^k(2^\delta - 1)) < 1/(2^\delta - 1)
  < 1/(\delta \ln 2)$.  Apply~\eqref{eq:delta-small-k}.
\end{proof}

\begin{proposition}[Pillai-type finiteness]\label{prop:pillai}
  For each fixed positive integer~$q$, the equation
  $(3^k - 1)/d(k) = q$ (equivalently,
  $(q+1) \cdot 3^k - q \cdot 2^S = 1$)
  has at most finitely many solutions in~$k$, with
  $k$ effectively bounded.
\end{proposition}

\begin{proof}
  Fix $q \geq 1$.  The equation $(q+1) \cdot 3^k = q \cdot 2^S + 1$
  gives
  \[
    \frac{(q+1) \cdot 3^k}{q \cdot 2^S} = 1 + \frac{1}{q \cdot 2^S},
  \]
  so $\Lambda_q := S \ln 2 - k \ln 3 - \ln((q+1)/q)$
  satisfies $|\Lambda_q| < 1/(q \cdot 2^S) < 2^{-S}$.
  Rewrite as $\Lambda_q = S \ln 2 - k \ln 3 - c_q$ where
  $c_q = \ln((q+1)/q)$ is a fixed constant.
  Since $S \ln 2 - k \ln 3 = \delta \ln 2$ and
  $\delta = c_q/\ln 2 + O(2^{-S})$, we have
  $\delta \approx c_q/\ln 2$ for solutions.
  Substituting into the LMN bound~\eqref{eq:lmn-full}
  (applied to $\Lambda_0 = S \ln 2 - k \ln 3$,
  which satisfies $|\Lambda_0 - c_q| < 2^{-S}$):
  either $|\Lambda_0| > \exp(-8174)$ (by \cref{cor:delta-bound}),
  in which case $|\Lambda_0 - c_q| < 2^{-S}$ forces
  $2^{-S} > |c_q| - e^{-8174}$, which fails for
  $S > \log_2(1/|c_q|) + 1$ (since $c_q > 0$);
  or $|\Lambda_0| < c_q + 2^{-S}$, which constrains $\delta$.

  More precisely: the equation forces $\delta = c_q/\ln 2 + O(2^{-S})$,
  but by the LMN bound, $|S\ln 2 - k \ln 3| > \exp(-8174)$.
  If $c_q > e^{-8174}$ (which holds for all $q < e^{8174}$),
  then $|S\ln 2 - k \ln 3 - c_q| < 2^{-S}$ combined with
  $|S\ln 2 - k \ln 3| > e^{-8174}$ constrains~$S$:
  by the three-distance theorem for the sequence
  $\{n \log_2 3\}_{n \geq 1}$, the values of $\delta$ near
  $c_q/\ln 2$ are separated by at least $e^{-8174}$, so
  at most $O(1)$ solutions exist per~$q$.

  This is a well-studied instance of the Pillai equation
  $a^x - b^y = c$
  (Pillai~1936, Mih\u{a}ilescu~\cite{Mihailescu2004});
  see also Bugeaud--Laurent--Mignotte--Siksek~\cite{BLMS2006}
  for refined effective bounds.
\end{proof}

\subsection{Toward resolving Hypothesis~H}\label{ssec:closing-H}

\begin{proposition}\label{prop:KH}
  There exists an effectively computable constant $K_H$ such
  that Hypothesis~H holds for all $k > K_H$.
  Specifically, for any $q \leq e^{8174}/(\ln 2)^2$ and
  any solution~$k$ of the Pillai equation
  $(q+1) \cdot 3^k - q \cdot 2^S = 1$: $k \leq K_H$.
\end{proposition}

\begin{proof}
  By \cref{prop:pillai}, each $q$ contributes at most
  finitely many solutions.  Since $q$ is bounded
  (\cref{lem:q-bound}), $K_H = \max_{q \leq Q_{\max}}
  k_{\max}(q)$ is finite.
\end{proof}

\begin{remark}\label{rem:KH-computation}
  An explicit computation of $K_H$ requires applying the
  Baker--W\"ustholz theorem~\cite{BakerWustholz1993}
  or the refined bounds of
  Bugeaud--Laurent--Mignotte--Siksek~\cite{BLMS2006}
  to the specific Pillai equation.
  A crude estimate using the general Baker--W\"ustholz constants
  for three logarithms gives $K_H$ potentially as large
  as~$10^{12}$, but the specialized structure of
  our equation (the coefficients $q+1$ and~$q$ differ by~$1$)
  and the use of LMN-type two-logarithm bounds
  (by fixing~$q$ and treating $\ln((q+1)/q)$ as a constant)
  should reduce~$K_H$ substantially.
  A detailed computation of~$K_H$ with optimized Baker
  constants is left to future work.

  In any case, Hypothesis~H is \emph{effectively decidable}:
  its resolution is a finite computation, not an open
  mathematical problem.  Our verification up to $k = 50000$
  is a first step; extending to $k = K_H$ would make
  \cref{thm:B} unconditional.
\end{remark}

\subsection{Proof of Theorem~B}\label{ssec:proof-B}

\begin{proof}[Proof of \cref{thm:B}]
  Assume Hypothesis~H.
  \begin{enumerate}[label=(\roman*)]
    \item $k = 3, 5$: enumeration (\cref{thm:finite}).
    \item $k = 4$: phantom; no cycle by~\cite{SimonsDeWeger2005}.
    \item $k = 6, \ldots, 10000$: Lean verification
          (\cref{thm:finite}).
    \item $k = 10001, \ldots, K_A$: Python verification
          confirms Range Exclusion (all pass).
    \item $k > K_A$: $\Delta < d$ by \cref{prop:cond-a};
          $d \nmid (3^k-1)$ by Hypothesis~H.
          Range Exclusion holds. \qedhere
  \end{enumerate}
\end{proof}

% ============================================================
\section{Lean~4 formalization}\label{sec:lean}
% ============================================================

The proof of \cref{thm:A} is formalized in Lean~4
(version~4.28.0, no Mathlib dependency).

\subsection{Structure}

\begin{itemize}
  \item \texttt{Basic.lean}: $S(k)$, $d(k)$, $\corrSum$,
    composition enumeration, \texttt{checkAvoidance}.
  \item \texttt{RangeExclusion.lean}: $\corrSum_{\min}$,
    $\corrSum_{\max}$, \texttt{checkRE} implementing
    \eqref{eq:re-condition}, dispatch function
    \texttt{checkOne}, batch theorems.
\end{itemize}

The dispatch function routes:
$k = 3, 5 \to$ enumeration;
$k = 4 \to$ true (phantom, Simons--de~Weger axiom);
$k \geq 6 \to$ Range Exclusion check.

Ten batch theorems (each covering 1000 values) are proved
by \texttt{native\_decide}, giving a combined certificate
for $k = 3, \ldots, 10000$.

\subsection{Axioms}

\begin{enumerate}
  \item \texttt{simons\_de\_weger}: no cycle with $k < 68$
    (from~\cite{SimonsDeWeger2005}).
  \item \texttt{baker\_lmn (k $\geq$ 10001)}: \texttt{checkRE k = true}
    (from~\cite{LMN1995}).  This axiom encodes: for
    $k \geq 10001$, both Condition~A ($\Delta < d$) and
    Condition~B ($d \nmid (3^k-1)$) hold.  A more principled
    splitting into separate axioms for the two conditions is
    possible but not implemented.
\end{enumerate}

\subsection{Trust base}

\texttt{native\_decide} compiles the check to native C code
and verifies the result against the Lean kernel.
Trust assumptions: Lean~4 kernel, Lean compiler, C compiler.
These are standard in the ITP community.

The function $S(k)$ is computed using
$\lceil k \cdot 15850 / 10000 \rceil$ (integer arithmetic,
no floating point).  The approximation
$15850/10000 = 1.5850$ slightly underestimates
$\log_2 3 \approx 1.58496$, so
$\lceil k \cdot 1.5850 \rceil \leq \lceil k \log_2 3 \rceil$
for all~$k$.  Equality is verified by a separate check
for all $k \leq 100000$: for each~$k$,
$2^{S_{\text{ceil}}} > 3^k \geq 2^{S_{\text{ceil}}-1}$.

\begin{table}[ht]
\centering
\caption{Lean~4 formalization summary.}
\label{tab:lean}
\begin{tabular}{ll}
\toprule
Metric & Value \\
\midrule
Lean version & 4.28.0 \\
Mathlib dependency & None \\
\texttt{sorry} count & 0 \\
Axioms & 2 (SdW~2005, LMN~1995) \\
Values verified (\texttt{native\_decide}) & 9997 \\
Build time & 337\,s (Apple M1 Pro, 16\,GB, macOS 15.3) \\
\bottomrule
\end{tabular}
\end{table}

% ============================================================
\section{Discussion}\label{sec:discussion}
% ============================================================

\subsection{What is proved unconditionally}\label{ssec:unconditional}

\Cref{thm:A} establishes $N_0(d(k)) = 0$ for
$k \leq 10000$, extending Simons--de~Weger's $k \leq 68$
by a factor of~$147$.  The Lean certificate provides
machine-checked assurance.

\subsection{What requires Hypothesis~H}\label{ssec:conditional}

\Cref{thm:B} (all~$k$) is conditional on~H.
As shown in \cref{ssec:closing-H}, Hypothesis~H is
\emph{effectively decidable}: it reduces to a finite
computation ($k \leq K_H \approx 10^9$).
This is not an open mathematical problem but a computational
challenge.

\subsection{The $k = 4$ phantom}\label{ssec:phantom}

The phantom at $k = 4$ is the unique value
$k \leq 50000$ with $N_0(d(k)) \neq 0$.
No other phantom has been found.

\subsection{Limitations}\label{ssec:limits}

\begin{enumerate}
  \item \textbf{Hypothesis~H} is the main gap.
    Closing it requires verifying $d(k) \nmid (3^k-1)$
    for $k \leq K_H \approx 10^9$ (feasible but not done).
  \item \textbf{Upper bound formula.}
    $\corrSum_{\max} = 3^k + 3^r - 2$ is verified for
    $k \leq 20$, not formally proved.
    This does not affect \cref{thm:A}
    (\cref{rem:upper-safe}).
  \item \textbf{Scope.}
    Only positive cycles are excluded.
    The full Collatz conjecture (no divergent trajectories)
    requires different methods
    (e.g., Tao~\cite{Tao2022}).
\end{enumerate}

\subsection{Open questions}\label{ssec:open}

\begin{enumerate}
  \item Resolve Hypothesis~H computationally
    (verify $d(k) \nmid (3^k-1)$ for $k \leq K_H$).
  \item Improve the Baker constants to reduce~$K_H$.
  \item Prove $\corrSum_{\max} = 3^k + 3^r - 2$ formally
    and formalize in Lean~4.
  \item Extend Range Exclusion to negative cycles or
    generalized $qx{+}r$ dynamics.
\end{enumerate}

% ============================================================
\section*{Acknowledgements}
% ============================================================

The author thanks the developers of Lean~4 for the formal
verification infrastructure.
Computational work: Apple~M1~Pro, 16\,GB.
No external funding.

\medskip\noindent\textbf{Data availability.}
The Lean~4 source code and Python verification scripts
are available at
\url{https://github.com/eric-music/Collatz-Junction-Theorem}.

% ============================================================
\begin{thebibliography}{99}

\bibitem{BakerWustholz1993}
A.~Baker and G.~W\"ustholz,
\emph{Logarithmic forms and group varieties},
J. reine angew. Math. \textbf{442} (1993), 19--62.

\bibitem{BLMS2006}
Y.~Bugeaud, M.~Laurent, M.~Mignotte, and S.~Siksek,
\emph{Classical and modular approaches to exponential
Diophantine equations. I. Fibonacci and Lucas perfect powers},
Ann. of Math. \textbf{163} (2006), 969--1018.

\bibitem{Barina2021}
D.~Barina,
\emph{Convergence verification of the Collatz problem},
J. Supercomput. \textbf{77} (2021), 2681--2688.

\bibitem{BoehmSontacchi1978}
C.~B\"ohm and G.~Sontacchi,
\emph{On the existence of cycles of given length in integer
sequences like $x_{n+1} = x_n/2$ if $x_n$ even, and
$x_{n+1} = 3x_n + 1$ otherwise},
Atti Accad. Naz. Lincei Rend. \textbf{64} (1978), 260--264.

\bibitem{Crandall1978}
R.~E.~Crandall,
\emph{On the ``$3x + 1$'' problem},
Math. Comp. \textbf{32} (1978), 1281--1292.

\bibitem{Eliahou1993}
S.~Eliahou,
\emph{The $3x + 1$ problem: new lower bounds on nontrivial
cycle lengths},
Discrete Math. \textbf{118} (1993), 45--56.

\bibitem{Lagarias1985}
J.~C.~Lagarias,
\emph{The $3x + 1$ problem and its generalizations},
Amer. Math. Monthly \textbf{92} (1985), 3--23.

\bibitem{Lagarias2010}
J.~C.~Lagarias (ed.),
\emph{The Ultimate Challenge: The $3x + 1$ Problem},
AMS, 2010.

\bibitem{LMN1995}
M.~Laurent, M.~Mignotte, and Y.~Nesterenko,
\emph{Formes lin\'eaires en deux logarithmes et d\'eterminants
d'interpolation},
J. Number Theory \textbf{55} (1995), 285--321.

\bibitem{Mihailescu2004}
P.~Mih\u{a}ilescu,
\emph{Primary cyclotomic units and a proof of Catalan's
conjecture},
J. reine angew. Math. \textbf{572} (2004), 167--195.

\bibitem{SimonsDeWeger2005}
D.~Simons and B.~de~Weger,
\emph{Theoretical and computational bounds for $m$-cycles of
the $3n + 1$ problem},
Acta Arith. \textbf{117} (2005), 51--70.

\bibitem{Steiner1977}
R.~P.~Steiner,
\emph{A theorem on the Syracuse problem},
Proc. 7th Manitoba Conf. Numer. Math. Comput.,
1977, 553--559.

\bibitem{Tao2022}
T.~Tao,
\emph{Almost all orbits of the Collatz map attain almost bounded
values},
Forum Math. Pi \textbf{10} (2022), e12.

\end{thebibliography}

\end{document}
